\chapter{Sparse Intervention Control}
\label{ch:sparse-control}

The previous chapter optimized survival on held-out chronics. This chapter adds
a second objective: sparsity. A policy that keeps the grid alive by
reconfiguring substations at every opportunity is impractical for a control
room. Interventions carry operational costs, and an agent acting constantly is
not trustworthy. The question is whether decentralized
agents can keep high survival while intervening rarely, locally, and only when
necessary.
\\
\\
In the literature, sparse activity is most often enforced by hard-coding a
physical threshold that decides when the agent is allowed to act. Instead of
relying on an external override, this chapter tries to make sparse intervention
something the agent learns by itself.
\\
\\
Mechanisms are therefore measured on two competing axes: \revision{mean survival} and
intervention sparsity. A good mechanism must improve sparsity without hurting
survival.


\paragraph{Question.}
Can the agents keep high survival while behaving more like operators: acting
rarely and mainly when the grid state makes an intervention
useful?

\paragraph{Baseline and change.}
This block starts from the flat decentralized MAPPO topology-control actor. For
agent \(i\), the baseline policy is one categorical distribution over the full
local action space:
\[
  \pi_i(a_i \mid o_i),
  \qquad
  a_i \in \{0,1,\ldots,|\mathcal A_i|-1\}.
\]
Action \(0\) is the no-op action. During training, actions are sampled from the
policy, while deterministic evaluation uses the greedy action by default:
\[
  a_{i,t}^{eval}=\arg\max_a \pi_i(a\mid o_{i,t}).
\]
This chapter tests five ways to reduce unnecessary topology
changes: evaluation-time rho overrides, an intervention-gated actor, fixed
non-idle action penalties, adaptive intervention budgets, and behavioral cloning from heuristics.
\\
\\
The intervention gate separates whether to act from which action to take,
following the hierarchical splits used in topology control~\cite{manczak2023hrl,
vandersar2023hmarl}. The fixed and adaptive penalties instead treat non-idle
interventions as a limited resource, following constrained and budgeted
reinforcement learning~\cite{achiam2017cpo,stooke2020pid,carrara2019budgeted}.

\section{Evaluation-Time Rho Heuristics}
\label{subsec:sparse-control-heuristics}

\paragraph{Implementation detail.}
The rho heuristic is an evaluation-only diagnostic. It does not change the
training objective and it does not prove that the policy learned sparse
behavior. The policy first proposes the deterministic action \(\bar a_{i,t}\),
then the evaluator may replace it with action \(0\).
\\
\\
The global version uses the pre-action maximum line loading
\[
  \rho_t^{global}=\max_{\ell\in\mathcal L}\rho_{\ell,t},
\]
and executes
\[
  a_{i,t}^{exec}=
  \begin{cases}
    0, & \rho_t^{global}<\rho_{safe},\\
    \bar a_{i,t}, & \text{otherwise,}
  \end{cases}
  \qquad \rho_{safe}=0.90.
\]
The local version is decentralized. Agent \(i\) computes the maximum rho on the
lines touching its regional substations,
\[
  \rho_{i,t}^{local}=\max_{\ell\in\mathcal L_i}\rho_{\ell,t},
\]
and only that agent is forced to no-op when
\(\rho_{i,t}^{local}<\rho_{safe}\). A boundary line can appear in the local
neighborhood of both adjacent agents, but no communication is required: both
agents can observe the line through their own physical neighborhood.

\paragraph{Interpretation.}
The heuristic is useful to ask what happens if safe-state interventions are
blocked, and if succesful, it can prove that the agents act too often.

\section{Intervention-Gated Actor}
\label{subsec:sparse-control-gate}

\paragraph{Implementation detail.}
The gated actor changes the actor factorization. Instead of one categorical
head over all actions, each agent has a gate head and a non-idle action head:
\[
  \pi_i^{gate}(g_i\mid o_i),\qquad g_i\in\{0,1\},
\]
where \(0\) means do nothing and \(1\) means intervene, and
\[
  \pi_i^{nonidle}(b_i\mid o_i),\qquad
  b_i\in\{1,\ldots,|\mathcal A_i|-1\}.
\]
During training, the gate is sampled first. If \(g_{i,t}=0\), then
\(a_{i,t}=0\). If \(g_{i,t}=1\), the final action is sampled from the non-idle
head. The PPO log-probability of the executed action is therefore
\[
  \log P(a_i=0)=\log P(g_i=0),
\]
and, for \(a_i>0\),
\[
  \log P(a_i)=\log P(g_i=1)+\log P(b_i=a_i\mid g_i=1).
\]

\noindent \latestrevision{Training always uses the hierarchical sampling
procedure above. At deterministic evaluation, the resulting policy can be
decoded in two ways. The \texttt{hierarchical\_greedy} rule first selects the
most likely gate decision and, if it chooses intervention, the most likely
non-idle action. The \texttt{final\_action\_map} rule instead selects the most
likely executed action using the final probabilities defined above. MAP is
therefore not a separate training mode; it only replaces sampling with a
deterministic decision. It can favor action \(0\) because the total intervention
probability is divided among several non-idle actions. The entropy variant used
by the gate-plus-budget diagnostic is specified in
Appendix~\ref{app:sparse-control-method-details}.}

\paragraph{Interpretation.}
The gate is learned and decentralized, but it can limit exploration more than a
reward penalty does. If it collapses early to no-op, the non-idle head is rarely
used and gets little learning signal.

\section{Fixed Sparse-Action Penalty}
\label{subsec:sparse-control-fixed-penalty}

\paragraph{Implementation detail.}
The fixed-penalty experiments compare several penalty strengths. In each run,
a reward penalty is applied to the final executed action of each agent:
\[
  r'_{i,t}=r_{i,t}
  -\lambda_{fixed}\mathbf{1}[a_{i,t}\neq0],
\]
with $\lambda_{fixed}\in\{0.0,0.001,0.003,0.01\}$.
\latestrevision{The complete 24-run actor, penalty, and seed grid is reported
in Appendix~\ref{app:sparse-control-method-details}.}

\paragraph{Interpretation.}
This mechanism changes only the training reward. It does not override actions during
evaluation and it does not block exploration during training. The policy can
still sample non-idle actions, but each such action must be useful enough to pay
its fixed cost.

\section{Adaptive Intervention Budget}
\label{subsec:sparse-control-aib}

{\color{black}
\paragraph{Implementation detail.}
The adaptive intervention budget does not use a fixed intervention penalty.
Instead, each agent has its own penalty coefficient \(\lambda_i\), which is
adjusted after every rollout. The local-safe version used in
Figure~\ref{fig:sparse-control-fulltest-tradeoff} first converts the maximum
line loading visible to agent \(i\) into a safety weight:
\[
  w_i^{local}(s_t)=
  \sigma\!\left(k\left(\rho_{safe}-\rho_{i,t}^{local}\right)\right).
\]
When the local grid is safely below \(\rho_{safe}=0.90\), this weight is high.
It decreases as the local lines approach or exceed the threshold. The action
cost is then
\[
  c_i(s_t,a_{i,t})=
  \mathbf{1}[a_{i,t}\neq0]w_i^{local}(s_t).
\]
Doing nothing always has zero cost. An intervention is expensive in a safe
state but remains relatively cheap when the local grid is stressed and an
action may be necessary. During the rollout, this cost is subtracted from the
original reward:
\[
  r'_{i,t}=r_{i,t}-\lambda_i c_i(s_t,a_{i,t}).
\]
A large value of \(\lambda_i\) therefore makes interventions more expensive,
whereas a small value leaves the policy freer to act.
\\
\\
At the end of the rollout, the average observed cost is
\[
  \hat C_i=\frac{1}{T}\sum_{t=1}^{T}c_i(s_t,a_{i,t}).
\]
This value is compared with the target budget \(d\), and the penalty
coefficient is updated as
\[
  \lambda_i\leftarrow
  \mathrm{clip}\left(
    \lambda_i+\eta(\hat C_i-d),
    0,
    \lambda_{\max}
  \right).
\]
If \(\hat C_i>d\), the agent has exceeded its budget, so \(\lambda_i\)
increases and future interventions become more expensive. If
\(\hat C_i<d\), the penalty decreases. The budget \(d\) is therefore a target
used to adjust the penalty after each rollout; it is not a fixed amount
subtracted from every reward. The value \(\rho_{safe}=0.90\) is only the center
of the smooth safety weight; it does not force the agent to do nothing. The
multiplier hyperparameters and complete budget ablation are reported in
Appendix~\ref{app:sparse-control-method-details}.
}

\section{Behavioral Cloning from Heuristics}
\label{subsec:sparse-control-bc}

\latestrevision{A teacher-student diagnostic tests whether the behavior of the
local rho heuristic can be copied into the policy weights. The student is
trained by behavioral cloning and evaluated without the heuristic. The
training setup and exact results are reported in
Appendix~\ref{app:sparse-control-method-details}.}

\section{Results}
Figure~\ref{fig:sparse-control-fulltest-tradeoff} summarizes \revision{mean survival} and the
mean action-0 fraction across agents. Exact condition means are reported in
Appendix~\ref{app:sparse-control-full-test-tables}.

\begin{table}[htbp]
  \centering
  \small
  \caption{Sparse-control mechanisms and where they affect the policy.}
  \label{tab:sparse-control-mechanisms}
  \setlength{\tabcolsep}{3pt}
  \begin{tabular}{L{0.18\textwidth}L{0.30\textwidth}L{0.24\textwidth}L{0.18\textwidth}}
    \toprule
    Method & Training effect & Evaluation effect & Main risk \\
    \midrule
    Baseline & Samples from the flat policy & Greedy argmax by default &
    \latestrevision{No explicit mechanism to favor sparse control} \\
    Rho heuristic & No training change & Hard no-op override in safe states &
    Evaluation is manually changed \\
    Intervention gate & Changes action factorization and sampling &
    Final-action MAP or hierarchical greedy & Gate collapse can starve non-idle exploration \\
    Fixed penalty & Fixed reward penalty on non-idle actions & No action override &
    Penalty coefficient must be tuned \\
    AIB local-safe & Adaptive state-dependent reward penalty & No action override &
    Depends on the local rho safety-weight design \\
    Behavioral cloning & Distills heuristic into actor weights & No action override &
    Struggles with extreme action-0 class imbalance \\
    \bottomrule
  \end{tabular}
\end{table}

% \begin{table}[htbp]
%   \centering
%   \small
%   \caption{Main sparse-control comparisons to report.}
%   \label{tab:sparse-control-main-comparisons}
%   \setlength{\tabcolsep}{4pt}
%   \begin{tabular}{L{0.45\textwidth}L{0.45\textwidth}}
%     \toprule
%     Question & Comparison \\
%     \midrule
%     Does a fixed sparse penalty already work? &
%     Baseline versus flat actor with $\lambda_{fixed}=0.010$ \\
%     Does the gate help beyond a sparse objective? &
%     Flat versus gated actor at $\lambda_{fixed}=0.010$ \\
%     Fixed versus adaptive coefficient? &
%     Fixed non-idle penalty versus adaptive non-idle cost at $d=0.20$ \\
%     Does local-safe weighting matter? &
%     Local-safe versus plain non-idle cost at $d=0.20$ \\
%     Does global-safe weighting matter? &
%     Global-safe versus local-safe cost at $d=0.20$ \\
%     Is \(\rho_{safe}=0.90\) special? &
%     Local-safe cost with $\rho_{safe}\in\{0.85,0.90,0.95\}$ \\
%     Is the heuristic only an evaluation trick? &
%     Rho heuristics versus learned fixed-penalty and AIB policies \\
%     \bottomrule
%   \end{tabular}
% \end{table}

\begin{figure}[H]
  \centering
  \includegraphics[width=0.96\linewidth]{figures/sparse_control_fulltest_tradeoff.png}
  \caption{Joint action-0/\revision{mean-survival} comparison of the selected sparse-control
  mechanisms on \texttt{bus14}. Large markers show condition means and faint
  markers the individual seeds. The plot uses checkpoint-aligned action-0
  rates and full-test \revision{mean survival}, so it shows the chapter's two objectives at
  once.}
  \label{fig:sparse-control-fulltest-tradeoff}
\end{figure}

\paragraph{Interpretation.}
\latestrevision{Local-safe AIB is the clear winner among the learned methods.
For \(d=0.10\) and \(d=0.20\), it raises mean survival from the baseline's
93.67\% to about 98.4\%, while reducing the mean non-idle action fraction from
56.6\% to about 2.1--2.2\%. It therefore improves both objectives at the same
time instead of trading survival for sparsity. This result suggests that the
adaptive local-safe cost helps each agent learn to reserve interventions for
states in which they are useful, while remaining fully decentralized.}

\section{Preliminary WCCI Stress Test}
\label{sec:sparse-wcci-stress-test}

\latestrevision{As an early scale test, representative mechanisms were
retrained from scratch with a fixed-width MLP actor on the
maintenance-enabled WCCI grid. The mechanisms changed intervention frequency,
but none reliably improved survival over the weak baseline, and the differences
were small relative to seed variation. Because this pilot changes the grid,
action space, and maintenance setting at the same time, it cannot isolate the
cause. It nevertheless motivates the size-independent graph actor introduced
next.}
\\
\\
\begin{tcolorbox}[
  colframe=red!75!black,
  colback=red!5!white,
  title=Chapter Summary
]
\newrevision{\textbf{Purpose:} Keep survival while reducing unnecessary
topology interventions.}
\\
\\
\newrevision{\textbf{Content:} Compare evaluation-time rho overrides, a
learned intervention gate, fixed action penalties, adaptive intervention
budgets, and behavioral cloning from a sparse teacher.}
\\
\\
\newrevision{\textbf{Findings:}}
\begin{itemize}
  \item \newrevision{\textbf{Evaluation heuristics} reach 97--98\% mean
  survival with action 0 in 94--97\% of decisions, but impose sparsity after
  the policy acts.}
  \item \newrevision{\textbf{A fixed penalty} already gives a strong learned
  controller: $\lambda_{fixed}=0.010$ reaches 97.44\% mean survival and a
  0.905 action-0 fraction without an evaluation override.}
  \item \newrevision{\textbf{Local-safe AIB} gives the best observed learned
  trade-off in this experiment: budgets $d=0.10$ and $d=0.20$ reach about
  98.4\% mean survival and 0.98 action-0 usage. The gated variants are less
  reliable and can collapse to almost permanent no-op behavior.}
\end{itemize}
\vspace{0.35cm}
\newrevision{\textbf{Takeaway:} Sparse behavior can be learned on
\texttt{bus14}, but the preliminary maintenance-enabled WCCI stress test does
not isolate transfer and its fixed-width MLP policies remain weak. This
motivates the graph representation introduced next.}
\end{tcolorbox}
